%% kapa: exam=1 round=36 year=2025 period=1 form=A subject=economics subject_ko=경제학원론 seq=29 number=69
%% source: Q-Net 감정평가사 기출문제 (cst00309) artlSeq=5231525 file=raw/1008/5231525/2025년 제36회 감정평가사 1차시험 1교시 문제지 원본.pdf
%% conversion: kordoc 4.24.51 -> tools/gen1.py (auto); review=proofread-sonnet; flags=none
\begin{question}{69}
자본이동이 완전한 소규모 개방경제의 먼델-플레밍모형하에서 변동환율제인 경우, 중앙은행이 통화량을 증가시킬 때 발생하는 단기적 효과에 관한 설명으로 옳은 것을 모두 고른 것은? (단, 국내외 국가의 물가는 고정)
\begin{framedtext}
\begin{tabular}{@{}p{0.45\linewidth}l@{}}
ㄱ. 순수출 증가 & ㄴ. IS곡선의 우측 이동 \\
ㄷ. LM곡선의 우측 이동 & ㄹ. 자국화폐의 평가절상
\end{tabular}
\end{framedtext}
\begin{choices}
  \item ㄱ, ㄴ
  \item ㄱ, ㄷ
  \item ㄴ, ㄷ
  \item ㄴ, ㄹ
  \item ㄷ, ㄹ
\end{choices}
\end{question}
\answerall
\answernote{정답표에 1,2,3,4,5 모두 정답으로 표기(전원 정답)}
